\chapter{Getting Started}
\label{ch:getting_started_pyclical}

\section{Installation}
PyClical is part of the \glucat project. To install it, you generally follow the standard \glucat installation procedure:

\begin{verbatim}
./configure
make
sudo make install
\end{verbatim}

This installs the PyClical Python extension module, typically into your site-packages directory.
For detailed installation instructions, including dependencies on Cython and Boost, refer to the \texttt{INSTALL.md} file in the source distribution.

\section{Running PyClical}
To use PyClical, start your Python interpreter (e.g., \texttt{python3} or \texttt{ipython3}) and import the module:

\begin{verbatim}
>>> from PyClical import *
\end{verbatim}

You can verify your installation by creating a simple multivector:

\begin{verbatim}
>>> e(1)
{1}
\end{verbatim}

\section{Demos and Tutorials}
The distribution includes a set of tutorials and demos in the \texttt{pyclical/demos} directory.
You can run the interactive tutorial menu by executing:

\begin{verbatim}
python3 pyclical_tutorials.py
\end{verbatim}
